\documentclass[11pt,a4paper]{article}

\usepackage[T1]{fontenc}
\usepackage[utf8]{inputenc}
\usepackage[spanish]{babel}
\usepackage[margin=2.3cm]{geometry}
\usepackage{graphicx}
\usepackage{subcaption}
\usepackage{amsmath,amssymb}
\usepackage{booktabs}
\usepackage[hidelinks]{hyperref}

\setlength{\parindent}{0pt}
\setlength{\parskip}{4pt}
\captionsetup{font=small,labelfont=bf}

\newcommand{\fig}[3]{%
  \IfFileExists{#1}{%
    \includegraphics[width=#2]{#1}%
  }{%
    \fbox{\parbox[c][2.2cm][c]{#2}{\centering\small [Figura no disponible: #3]}}%
  }%
}

\title{\textbf{Preparación de datos para estudiar la relación entre ENSO y SWE}\\
\large Pacific Northwest}
\author{Nicolás Mínguez Aguirre}
\date{Curso académico 2026--2027}

\begin{document}

\maketitle

\begin{abstract}
Este trabajo prepara un conjunto de datos para estudiar si la sensibilidad del
equivalente en agua de la nieve (\emph{Snow Water Equivalent}, SWE) al índice
ENSO (\emph{El Niño--Southern Oscillation}) cambia con la elevación y la zona
geográfica del Pacific Northwest (PNW). El dataset utiliza una fila por
estación y temporada, combina observaciones mensuales de estaciones SNOTEL con
información geográfica e índice ENSO, y conserva indicadores de completitud.
El análisis se divide en tres objetivos: identificar los puntos existentes,
determinar cuáles son comparables y utilizables, y describir qué relación puede
observarse en los datos tratados. La contribución principal es la preparación
reproducible del dataset; el modelo final se mantiene deliberadamente sencillo.
\end{abstract}

\section{Introducción}

El equivalente en agua de la nieve (SWE) representa la cantidad de agua
contenida en un manto nivoso. En el Pacific Northwest (PNW), la nieve
acumulada durante el invierno constituye una reserva hídrica importante para la
primavera y el verano, cuando contribuye al abastecimiento, al riego y a los
caudales de los ríos \cite{she2007}. La misma reserva puede relacionarse con
riesgos opuestos: una acumulación o una fusión rápida puede contribuir a
inundaciones, mientras que una acumulación baja puede reducir el agua
disponible y aumentar el riesgo de sequía.

La temperatura y la precipitación ayudan a explicar el SWE. Una temperatura
más alta puede hacer que parte de la precipitación caiga como lluvia en vez de
nieve, y la cantidad de precipitación determina cuánto material puede
acumularse. Sin embargo, estas variables no describen por sí solas todo el
proceso. El SWE también depende de la circulación atmosférica, el viento, la
humedad, la fecha en que cae la precipitación y las interacciones entre
temperatura y precipitación. Además, puede existir una tendencia a largo plazo
relacionada con el calentamiento u otros cambios atmosféricos. Algunos de estos
factores no se observan en el dataset o no se incluyen en el alcance de esta
práctica.

\begin{figure}[ht]
  \centering
  \begin{subfigure}{0.31\linewidth}
    \centering
    \fig{images/tavg_2015.jpg}{\linewidth}{temperatura de 2015}
    \caption{Temperatura}
  \end{subfigure}
  \hfill
  \begin{subfigure}{0.31\linewidth}
    \centering
    \fig{images/prcp_2015.jpg}{\linewidth}{precipitación de 2015}
    \caption{Precipitación}
  \end{subfigure}
  \hfill
  \begin{subfigure}{0.31\linewidth}
    \centering
    \fig{images/swe_2015.jpg}{\linewidth}{SWE de 2015}
    \caption{SWE}
  \end{subfigure}
  \caption{Contexto del PNW durante 2015. Las imágenes muestran que un año con
  condiciones térmicas y pluviométricas particulares puede producir un SWE
  excepcionalmente bajo.}
  \label{fig:2015}
\end{figure}

ENSO es una variable externa a la estación de nieve y puede proporcionar
información con antelación sobre condiciones estacionales. Estudios sobre el
oeste de Estados Unidos han relacionado la variabilidad oceánico-atmosférica
con la nieve \cite{hunter2006}. El interés de este trabajo no es afirmar que
ENSO explique todo el SWE, sino estudiar qué parte de la variación se asocia
con ENSO y si esa asociación cambia según la estación.

Como primera ilustración se considera Stampede Pass, una estación de elevación
relativamente alta. En la serie temporal de la estación se observan tres
características. Primero, aparece una tendencia general descendente. Esta
tendencia podría estar relacionada con el calentamiento a largo plazo, cambios
en la circulación atmosférica u otros factores no observados. Si no se
considera, parte de su variación podría atribuirse incorrectamente a ENSO.
Segundo, existe una incertidumbre considerable: años con valores parecidos de
ENSO pueden tener SWE muy diferente. Tercero, se aprecia una tendencia
descriptiva a que los años de La Niña tengan más SWE que los años de El Niño,
aunque no se cumple en todos los años.

\begin{figure}[ht]
  \centering
  \begin{subfigure}{0.48\linewidth}
    \centering
    \fig{images/year_enso_vs_swe.png}{\linewidth}{serie temporal}
    \caption{SWE por temporada}
  \end{subfigure}
  \hfill
  \begin{subfigure}{0.48\linewidth}
    \centering
    \fig{images/enso_vs_swe.png}{\linewidth}{relación ENSO--SWE}
    \caption{ENSO y SWE}
  \end{subfigure}
  \caption{Ejemplo de una estación del dataset. Cada punto representa una
  temporada.}
  \label{fig:stampede}
\end{figure}

La segunda figura muestra una relación aparente entre ENSO y SWE en esta
estación concreta. Esa relación debe interpretarse teniendo en cuenta su zona
y su elevación: no sabemos todavía si la misma pendiente aparecería en otras
estaciones. La pregunta central es, por tanto, si el efecto de ENSO cambia
entre zonas y elevaciones.

Conceptualmente, una primera comparación consiste en representar SWE frente a
ENSO y ajustar una línea para cada estación. Las líneas pueden compararse
posteriormente con la elevación y la zona. La hipótesis es que la sensibilidad
al ENSO no es uniforme. Las estaciones cercanas al límite lluvia--nieve podrían
mostrar una respuesta más fuerte porque pequeños cambios térmicos alteran el
tipo de precipitación. Por otro lado, una zona como Mount Baker puede recibir
más nieve que otras zonas incluso a elevaciones similares. El nivel medio de
SWE y su sensibilidad al ENSO son, por tanto, cuestiones relacionadas pero
distintas.

El trabajo se organiza en tres objetivos:

\begin{enumerate}
  \item \textbf{Objetivo 1: ¿qué puntos existen?} Construir el dataset inicial.
  \item \textbf{Objetivo 2: ¿qué puntos podemos usar?} Tratar la completitud y
  comparabilidad de estaciones y temporadas.
  \item \textbf{Objetivo 3: ¿qué indican los puntos utilizables?} Proponer un
  uso descriptivo y un modelo mínimo.
\end{enumerate}

\section{Objetivo 1: ¿qué puntos existen?}

\subsection{Fuentes y estaciones}

El dataset se construye usando una única fuente principal de datos terrestres:
la red SNOTEL consultada mediante AWDB. SNOTEL es la red de estaciones de
observación de nieve utilizada en el proyecto para recoger información de
estaciones de montaña. Limitar la selección a estaciones SNOTEL evita mezclar
redes con definiciones, instrumentos o procedimientos de observación
potencialmente distintos.

Se seleccionan las estaciones de la región hidrográfica configurada en el
proyecto y dentro del intervalo de elevación de 500--4500 m. Para cada estación
se conservan identificadores, nombre, latitud, longitud, elevación y zona
geográfica. La clasificación espacial utiliza las unidades HUC disponibles en
el repositorio y una tabla de nombres de zonas.

\subsection{Temporadas y variables}

Se descargan medidas mensuales de temperatura media, precipitación y SWE para
los meses de noviembre a marzo. El SWE suele comunicarse mediante la medida del
1 de abril, cuando normalmente se encuentra cerca del máximo de la reserva
nival \cite{she2007}. En este proyecto se utiliza el máximo de SWE observado en
la temporada. Esta elección evita depender estrictamente de una fecha fija si
la acumulación máxima ocurre antes o después del 1 de abril.

Los datos mensuales se agregan a una fila por estación y temporada:

\begin{itemize}
  \item temperatura: media de los meses disponibles;
  \item precipitación: suma de los meses disponibles;
  \item SWE: máximo mensual de la temporada;
  \item ENSO: media del índice mensual durante la temporada.
\end{itemize}

También se guardan las variables \(n_{\mathrm{TAVG}}\),
\(n_{\mathrm{PRCP}}\), \(n_{\mathrm{SWE}}\) y \(n_{\mathrm{ENSO}}\), que indican
cuántos meses contribuyeron a cada resumen. Un dato ausente se conserva como
ausente y no se convierte automáticamente en cero.

Las observaciones mensuales consecutivas no son completamente independientes.
La nieve, la temperatura y la precipitación pueden persistir entre meses, y el
índice ENSO también tiene una estructura temporal. Agregar los meses a
temporadas reduce la dependencia mes a mes y produce una unidad más cercana a
la pregunta científica. El índice ENSO se agrega de la misma forma para que
cada temporada de una estación pueda asociarse con un único valor ENSO.

El resultado de Objetivo 1 es un CSV con una fila por estación--temporada. Este
archivo contiene los puntos disponibles, pero todavía puede incluir temporadas
incompletas, estaciones con pocos años y distribuciones espaciales desiguales.
Por ello, la recopilación no es todavía una base lista para el análisis.

\section{Objetivo 2: ¿qué puntos podemos usar?}

\subsection{Completitud y comparabilidad}

La primera decisión es eliminar las estación--temporada que no tengan cinco
mediciones mensuales de SWE:

\[
n_{\mathrm{SWE}}=5.
\]

Un máximo calculado a partir de cuatro meses no es directamente comparable con
un máximo calculado a partir de cinco. Además, el efecto de la ausencia depende
del mes que falte: no es equivalente perder noviembre que perder el mes en que
normalmente ocurre el máximo.

No se imputan los valores de SWE faltantes. Imputar a partir de otras
estaciones, elevaciones o zonas podría parecer razonable, pero también podría
introducir precisamente la estructura ENSO--SWE que se intenta estudiar. La
imputación podría aumentar o reducir artificialmente la pendiente, disminuir
aparentemente la incertidumbre y hacer que la relación dependa del modelo de
imputación. Para este proyecto es más transparente excluir las temporadas
incompletas y conservar los valores originales.

Un cero de SWE no se trata como un valor ausente. Si el cero procede de una
observación válida, se conserva; solamente se eliminan las temporadas sin las
cinco mediciones requeridas.

\subsection{Cobertura temporal y control de calidad}

Después de conservar las temporadas completas, se mantienen las estaciones con
al menos 30 temporadas utilizables. Este umbral no es una propiedad universal
de los datos, sino un criterio operativo: una línea ajustada con muy pocos años
tendría poco apoyo temporal y sería especialmente sensible a un solo año
extremo. Con 30 temporadas se conserva una parte sustancial del registro
disponible sin exigir una cobertura perfecta.

Los valores extremos no se eliminan automáticamente. Un SWE muy alto puede
representar una temporada real y no un error. Se puede marcar un valor como
candidato a revisión mediante un gráfico de densidad, un boxplot o una regla
del rango intercuartílico, pero solo se elimina si existe una razón
independiente: duplicación, unidad incorrecta, valor imposible o error
documentado en la fuente.

\subsection{Distribuciones y representatividad}

El dataset debe auditarse antes de modelar. Es necesario comprobar la
distribución de SWE, ENSO, elevación y zona:

\begin{itemize}
  \item SWE suele ser positivo y asimétrico; una densidad puede mostrar si unos
  pocos años de mucha nieve dominan la distribución;
  \item ENSO debe cubrir valores positivos y negativos, sin que se creen años
  artificiales para equilibrarlo;
  \item la elevación debe cubrir un rango útil para comparar estaciones;
  \item las zonas pueden contener números muy diferentes de estaciones.
\end{itemize}

\begin{figure}[ht]
  \centering
  \fig{images/inference_enso_swe.png}{0.70\linewidth}{distribuciones de las variables}
  \caption{Representación disponible del conjunto tratado. El gráfico muestra
  SWE frente a ENSO y utiliza el color para representar la elevación.}
  \label{fig:distribution}
\end{figure}

\textbf{[El repositorio no contiene un gráfico de densidad específico de SWE,
ENSO, elevación y zona; debe generarse a partir del CSV tratado.]}

No se pretende forzar distribuciones uniformes. Si una zona tiene pocas
estaciones, se informa de ello y sus comparaciones se interpretan con menos
confianza. La elevación se conserva como variable continua, en vez de crear
bandas artificiales. ENSO no se rebalancea ni se sobremuestrea porque es una
variable temporal común a todas las estaciones de un mismo año.

El resultado de Objetivo 2 es un CSV limpio en el que todas las estación--temporada
usadas tienen cinco mediciones de SWE y cada estación tiene al menos 30
temporadas utilizables. La correlación restante entre temporadas, estaciones y
zonas no se resuelve mediante la limpieza; se considera al utilizar los datos
en Objetivo 3.

\section{Objetivo 3: ¿qué indican los puntos?}

Este objetivo se mantiene breve porque el valor principal del trabajo es el
dataset. El uso propuesto es un modelo mínimo para describir la sensibilidad del
SWE al ENSO y comprobar si cambia con la elevación y la zona.

\subsection{Efecto temporal y dependencia}

La serie de Stampede Pass sugiere una tendencia temporal descendente. Para no
atribuir esa tendencia a ENSO, se incluye el año como una covariable lineal
centrada. El año no se divide en un efecto libre para cada temporada, porque
eso consumiría información y podría absorber la misma variación anual necesaria
para estudiar ENSO.

Las temporadas consecutivas pueden seguir relacionadas después de agregar los
meses. No se añade inicialmente una variable \(\Delta\mathrm{ENSO}\): cambiaría
la pregunta desde el efecto del ENSO de la temporada hacia una diferencia entre
temporadas. Si se necesitan intervalos de incertidumbre, la alternativa mínima
es remuestrear bloques de años completos, manteniendo juntas todas las
estaciones de cada año. Así se conserva el ENSO común y parte de la dependencia
temporal. La correlación se trata aquí como una limitación de la incertidumbre,
no como una razón para eliminar años válidos.

\subsection{Elevación y zona}

La elevación se incluye como variable continua. Dividirla en bandas podría
agrupar estaciones que experimentan efectos distintos o producir fronteras
artificiales. Se utiliza la versión centrada y escalada:

\[
E'_j =
\frac{E_j-\overline{E}}{\operatorname{SD}(E)}.
\]

La interacción ENSO--elevación es el término que responde a la hipótesis
principal: permite que la pendiente de SWE frente a ENSO cambie a lo largo del
gradiente altitudinal.

La zona se trata como una variable cualitativa:

\[
\log(\mathrm{SWE}_{jt}) =
\alpha_j+\beta_X X_t+\beta_E E'_j+
\beta_{XE}X_tE'_j+\gamma_{Z(j)}+\lambda T_t+\varepsilon_{jt}.
\]

Si los diagnósticos de la distribución y de los residuos lo justifican, se
puede usar \(\log(\mathrm{SWE})\), porque SWE es positivo y marcadamente
asimétrico. La transformación expresa cambios relativos y reduce el peso de
los valores extremadamente altos. Si existieran ceros válidos, no se aplicaría
el logaritmo sin una decisión adicional.

La interacción ENSO--zona puede añadirse para permitir una pendiente ENSO
diferente en cada zona. La interacción completa
ENSO--elevación--zona permitiría que la relación entre ENSO y elevación cambiase
según la zona. Con 20 zonas, esto equivale a muchas pendientes y términos
adicionales; por ello se considera una extensión, no el modelo principal.

No se modela una correlación espacial entre zonas. Las zonas no están ordenadas
como los años y su vecindad no garantiza que tengan el mismo régimen climático.
Además, las zonas están definidas geográficamente y pueden contener estaciones
con elevaciones muy diferentes. En un proyecto centrado en preparación de
datos, es más transparente tratar la zona como categoría y documentar su
representación que imponer una estructura espacial compleja.

\subsection{Resumen de decisiones}

La ruta final es:

\begin{enumerate}
  \item construir una fila por estación y temporada;
  \item conservar el máximo estacional de SWE y la media estacional de ENSO;
  \item eliminar estación--temporada con menos de cinco mediciones de SWE;
  \item no imputar SWE;
  \item conservar estaciones con al menos 30 temporadas completas;
  \item auditar valores extremos y distribuciones sin eliminar observaciones
  válidas automáticamente;
  \item usar el año para controlar la tendencia general;
  \item mantener la elevación continua y centrada;
  \item incluir la interacción ENSO--elevación;
  \item tratar la zona como categoría, dejando sus interacciones más complejas
  como extensión;
  \item reconocer la dependencia temporal y utilizar bloques de años si se
  calculan intervalos.
\end{enumerate}

El resultado no debe presentarse como una inferencia climática definitiva. Es una
preparación reproducible de datos acompañada de un primer modelo interpretable.
La pregunta final queda formulada de manera concreta: después de tratar las
temporadas incompletas, las diferencias persistentes entre estaciones y la
tendencia temporal, ¿cambia la asociación entre ENSO y SWE con la elevación y
entre zonas?

\begin{thebibliography}{9}

\bibitem{hunter2006}
Hunter, T., Tootle, G. y Piechota, T. (2006).
Oceanic-atmospheric variability and western U.S. snowfall.
\emph{Geophysical Research Letters}, 33, L13706.
doi:10.1029/2006GL026600.

\bibitem{she2007}
She, N. y Basketfield, D. (2007).
Exploratory analysis of snowpack in Washington State.
En \emph{World Environmental and Water Resources Congress 2007}.
ASCE. doi:10.1061/40927(243)80.

\bibitem{uw2024}
Office of the Washington State Climatologist (2024).
\emph{ENSO's Effects on Snow Water Equivalent over the Years}.
\url{https://climate.uw.edu/2024/01/09/ensos-effects-on-snow-water-equivalent-over-the-years/}

\bibitem{psl}
NOAA Physical Sciences Laboratory.
\emph{Niño 3.4 SST Index from the HadISST1.1}.
\url{https://psl.noaa.gov/data/timeseries/month/Nino34/}

\bibitem{rayner2003}
Rayner, N. A. et al. (2003).
Global analyses of sea surface temperature, sea ice, and night marine air
temperature since the late nineteenth century.
\emph{Journal of Geophysical Research}, 108(D14), 4407.

\bibitem{gelman2007}
Gelman, A. y Hill, J. (2007).
\emph{Data Analysis Using Regression and Multilevel/Hierarchical Models}.
Cambridge University Press.

\end{thebibliography}

\end{document}